\begin{pastexam}[label={px:ntut-109-ci-02}]{北科大資工所 109 學年度 第 1 題（節錄 (4)(5) 小題）}
（簡答）

Please answer the following questions shortly and concisely. It is assumed that $T(1)=d$ for some
constant $d$. (15 pts)
\begin{enumerate}[label=(\arabic*)]
\setcounter{enumi}{3}
\item State, using the ``big oh'' notation, the solution to $T(n)=8T\left(\frac{n}{2}\right)+n^3$.
(3 pts)
\item State, using the ``big oh'' notation, the solution to
\[T(n)=14T\left(\tfrac{n}{2}\right)+18n^3\sqrt{n}.\ \text{(3 pts)}\]
\end{enumerate}
\end{pastexam}

\begin{pastsoln}
答案：(4) $O(n^3 \log n)$；(5) $O(n^{\log_2 14})$，約 $O(n^{3.81})$。

兩題都是 $T(n) = aT(n/b) + f(n)$，$T(1) = d$ 是常數，套資結書定理~2.4 的主定理：先算
$n^{\log_b a}$，再和 $f(n)$ 比大小。

\textbf{(4)} $a = 8$、$b = 2$，$n^{\log_2 8} = n^3$。 $f(n) = n^3 = \Theta(n^3)$，同階，是情況 2：
$T(n) = \Theta(n^3 \lg n)$，所以 $T(n) = O(n^3 \log n)$。 用遞迴樹看（資結書第~2.3~節）：第 $i$
層有 $8^i$ 個子問題、每個大小 $n/2^i$，這一層的工作是 $8^i (n/2^i)^3 = n^3$，每一層一樣多，
共 $\lg n + 1$ 層，所以是 $n^3 \lg n$。 只寫 $O(n^3)$ 少了 $\lg n$，是錯的。

\textbf{(5)} $a = 14$、$b = 2$，$\log_2 14 \approx 3.807$。 $f(n) = 18 n^3 \sqrt{n} = 18 n^{3.5}$。 取
$\varepsilon = 0.3$，$\log_2 14 - 0.3 \approx 3.507 \ge 3.5$，所以 $f(n) = O(n^{\log_2 14 - \varepsilon})$，
$f$ 多項式地比 $n^{\log_2 14}$ 小，是情況 1：
\[
T(n) = \Theta(n^{\log_2 14}) \approx \Theta(n^{3.81}) ,
\]
所以 $T(n) = O(n^{\log_2 14})$。 常見的錯是看到 $n^{3.5}$ 比 $n^3$ 大就答 $O(n^{3.5})$；要比的是
$n^{\log_b a} = n^{3.807}$，不是 $n^3$。 常數 $18$ 不影響結果。
\end{pastsoln}
